\chapter{Performance Methodology and Benchmarking}
\label{ch:perf}

\section{What We Measure and Why}
\label{sec:perf-what}

Because MATHLIB5's correctness claims are independent of runtime, its
performance is reported only to characterize the implementation, never to
justify safety. This chapter states the methodology so that the figures in
\cref{tab:results} are reproducible and not over-interpreted.

\section{Harness}
\label{sec:perf-harness}

The benchmark harness:
\begin{enumerate}
  \item pins the input size $n$ (vector length) and seeds the APL source,
  \item runs the full ten-layer chain once for warm-up, discarding the
        result,
  \item runs $K=30$ repetitions, recording wall-clock per layer via
        receipt timestamps,
  \item reports the median and the inter-quartile range.
\end{enumerate}
The orchestration glue (Layer~0 scaffolding, \cref{ch:build}) is excluded
from the per-layer timing to avoid measuring the build system itself.

\section{Scaling Behavior}
\label{sec:perf-scaling}

The summation reduction is $O(n)$ in the APL and IR layers and $O(n)$ in
the emitted C$--$ fold; the Lean closed-form check is $O(1)$ once the
theorem is in the context. We therefore expect total verified time to grow
linearly with $n$ for the numeric layers and stay flat for the proof layer.
The matrix contraction adds an $O(mn)$ term, as expected.

\section{Overhead of Verification}
\label{sec:perf-overhead}

The dominant non-numeric cost is the Lean kernel check of the closed form
and the Ed25519 signing at each boundary. Both are constant per artifact
and small relative to the numeric loop for large $n$. For very small $n$
the verification overhead can exceed the numeric work; this is acceptable
because the value of the system is certification, not speed, on tiny
inputs.

\section{Reporting Discipline}
\label{sec:perf-disc}

We follow three rules when publishing numbers:
\begin{itemize}
  \item never present a median without its spread,
  \item never attribute a speedup to verification (verification is a cost,
        bought voluntarily),
  \item always state the toolchain commit so the run is reproducible
        against the anchor (\cref{ch:anchors}).
\end{itemize}
These rules keep the performance section honest and prevent the common
temptation to market verification as free.
